\section{Diskusjon og videre arbeid}
\label{sec:diskusjon}
\subsection{Diskusjon}

\subsubsection{Tap i drivverket og betydningen for simuleringen}
Utsnurringen viser at tapene er mye større enn antatt i konseptevalueringen. Det konstante tapsmomentet er omtrent ti ganger det simuleringen brukte, og opptrekk med snor gir omtrent \SI{2700}{o/min} i stedet for \SI{4000}{o/min}. Til sammen reduserer det simulert kjøredistanse fra rundt \SI{35}{m} til rundt \SI{3}{m} (figur~\ref{fig:sim-maalte-tap}). Rangeringen av konseptene i seksjon~\ref{sec:konseptevaluering} bygger på samme tapsmodell for alle tre, og endres trolig ikke. Lagringen i kulelagre og tannhjulsinngrepet er fortsatt bedre enn en O-ring på kjegler. De absolutte distansene må derimot regnes på nytt.

Med større tap lønner det seg å ta ut energien raskere, og beste utveksling flytter seg fra 5--6 til omtrent 3. Kurven er flat rundt optimum, og $i = 4{,}29$ gir \SI{8}{\percent} kortere distanse enn $i = 3$. Utvekslingen bør likevel ikke endres før tapene er målt på nytt, av tre grunner. Svinghjullagrene alene ser ut til å stå for omtrent \SI{1,1}{mN.m} av tapet, og en stor del av det er trolig fett som kan vaskes ut. Spredningen med tannhjulet frakoblet viser at sideveis gnissing bidrar, og den skal avstandshylsene fjerne. Og $c_2$ er bare målt opp til omtrent \SI{3100}{o/min}. Rekkefølgen videre er derfor å montere avstandshylsene, gjenta serien med tannhjulet frakoblet, vaske ut fettet, måle igjen, og deretter vurdere utvekslingen og trommeldiameteren på nytt.

Turtallet etter opptrekk er like viktig som tapene. Energien går med kvadratet av turtallet, så \SI{2700}{o/min} gir bare \SI{46}{\percent} av energien ved \SI{4000}{o/min}. Opptrekkene er gjort med TR-1, den minste trommelen. Den krever mest kraft mot slutten av trekket, og TR-2 og TR-3 skal vise om en større trommel gir høyere turtall for hånd.
\begin{comment}
\subsubsection{Vurdering opp mot problemstilling og produktkrav}
\subsubsection{Betydningen av konsept- og designvalg}
\subsubsection{Testenes pålitelighet og begrensninger}
\subsubsection{Produktutviklingsprosess og forbedringsmuligheter}
% Skill mellom observasjoner, tolkninger og hypoteser som krever videre undersøkelser.
\end{comment}
